\documentclass[12pt,letterpaper]{article}

\usepackage[margin=1in]{geometry}
\usepackage[T1]{fontenc}
\usepackage{amsmath}
\usepackage{amssymb}
\usepackage{newtxtext,newtxmath}
\usepackage{setspace}
\usepackage{caption}
\usepackage{graphicx}
\usepackage{titlesec}
\usepackage[hidelinks]{hyperref}

% Match the blueprint's body-size, bold, unnumbered section headings.
\titleformat{\section}{\normalfont\normalsize\bfseries}{}{0pt}{}
\titlespacing*{\section}{0pt}{\baselineskip}{0pt}
\setlength{\parindent}{0pt}
\setlength{\parskip}{0pt}

\begin{document}

\begin{titlepage}
    \centering
    \singlespacing
    \vspace*{1.5in}

    {\bfseries Does higher education complement or substitute for family background?\par}
    \vspace{\baselineskip}
    Econ 4331W Development Economics\par
    November/11/2026\par
    Ellie Hou\par

    \vfill
\end{titlepage}

\doublespacing
\raggedright

\section*{INTRODUCTION}
Does completing a four-year college degree give students from a disadvantaged background a
larger earnings premium? 


This question helps inform a policy trade-off: to help students from
disadvantaged backgrounds thrive, do we simply lower tuition, or is a reform of the whole
education system needed? With data from the National Longitudinal Survey of Youth 1997,
this paper examines the difference in the college earnings premium by family background,
controlling for pre-college academic preparation, measured by ASVAB scores, and other
demographic characteristics.


\section*{LITERATURE REVIEW}

Mincer (1974) links earnings to schooling and work experience, treating these as human
capital investments in earning capacity. The framework provides a foundation for answering
my main research question by measuring how schooling can affect earning capacity.
Furthermore, Becker and Tomes (1986) asked a broader question: how do parents influence
their children’s human capital? In their model, parents prefer to optimize the amount of
investment in human capital before starting to save for voluntary bequests.
These questions also connect to the literature on wealth concentration and inequality. De
Nardi (2015) reviews basic Bewley-style models, in which households save to protect
consumption against earnings risk, alongside models incorporating bequests and
intergenerational human-capital transmission. The model helps explain the role of family investment in
differences in earning capacity. In this paper, I will focus on a narrower question: “Is college
completion associated with narrowing income disparities among different family-background
groups?”
More directly, Hendricks, Koreshkova, and Leukhina (2026) study the potential effect of
widening access to selective public universities on earning capacity. Their results show that
expanding enrollment capacity can increase the population’s aggregate earning capacity,
partly by admitting academically well-prepared students who were previously excluded by
admission constraints. My analysis, on the other hand, examines whether the adjusted
college earnings premium varies with family background.

\section*{HYPOTHESIS}
% Text carried over from the submitted preliminary report for revision.
My hypothesis is that students from lower-income families have a higher adjusted college
earnings premium; that is, $\delta>0$ in the model described in the methodology section.
The proposed mechanism behind my hypothesis is that college education can substitute
for family background and thus narrow differences in human capital among students from
different backgrounds.

% TODO: Support the proposed mechanism for the secondary hypothesis with literature.
My secondary hypothesis is that students from lower-income families are less likely to
complete a bachelor's degree. The reasoning behind this hypothesis is that, if schools
value outcomes associated with future earnings, colleges are incentivized to invest more
in students with higher demonstrated potential, and that potential can logically be
highly correlated with students' backgrounds.

\section*{DATA}
\textbf{Data source and sample.}
For this paper, I used the data from NLSY97 (the National Longitudinal Survey of Youth)
to construct variables. The survey includes a representative sample of Americans born
between 1980 and 1984. The source variables I included are household income in 1996,
parent-reported household net worth in 1997, wage income, and highest degree earned in 2015.

% TODO: Audit earnings receipt filters before revising the sample. Some valid skips
% can represent no wage income, and the current extract omits the receipt questions.
The balanced panel includes 2,123 people with all required variables observed in three
earnings years, giving 6,369 person-year observations. All respondents have a positive
custom weight supplied for this selected panel. Missing responses and valid skips are
excluded; reported zero earnings and genuine negative baseline income and net worth
are retained.

\textbf{Dependent variable $Y_{it}$.}
$Y_{it}$ is the annual wage and salary income of the surveyed person in 2016, 2018,
and 2020, measured in 2025 US dollars. These are the earnings years available in the
extract, collected in the 2017, 2019, and 2021 surveys.

\textbf{Independent variables $C_i$ and $L_i$.}
% TODO: Justify pooling degree categories. No bachelor's degree does not mean
% never attended college; degree status is fixed at the 2015 interview.
$C_i$ indicates whether the person completed a bachelor's degree or higher by the
2015 survey. $L_i$ indicates whether their baseline household income and net worth
are both below their respective 60th percentiles. The weighted cutoffs are \$55,000
for 1996 income and \$101,000 for 1997 net worth, in nominal dollars. I calculate each
cutoff within the selected balanced panel using
\href{https://www.nlsinfo.org/content/weights/nlsy97}{custom NLSY97 weights}, counting
each person once. Each cutoff is the smallest observed value at which cumulative
weight reaches 60\%. The lower-background group contains 1,151 people and represents
48.0\% of the weighted panel because both conditions must hold.
% TODO: Motivate the joint income-and-wealth definition and report sensitivity
% to the cutoffs already explored; do not select the definition by significance.

\textbf{Control variables $X_{it}$.}
I control for pre-college academic achievement using the NLSY97-created ASVAB
math--verbal percentile score, which is normalized within three-month birth groups.
In addition, I control for sex, age at the end of the earnings year, and race/ethnicity.
% Verified in the archived codebook: R9829600 is the age-normalized percentile;
% dividing by 1,000 restores its three implied decimal places.

% TODO: Choose the summary-table presentation after revising the sample.
\textit{[Insert your summary statistics table here.]}

\section*{METHODOLOGY}
I convert nominal dollars into constant 2025 dollars using annual average CPI-U,
multiplying each value by $CPI_{2025}/CPI_t$.\footnote{The supplied CPIAUCNS series
has no October 2025 observation, consistent with
\href{https://www.bls.gov/news.release/archives/cpi_12182025.htm}{BLS documentation}.
I use the mean of the 11 available months, 321.943.}

The primary model is
\[
    Y_{it}=\alpha+\beta C_i+\gamma L_i+\delta(C_i\times L_i)
           +X'_{it}\theta+\lambda_t+\epsilon_{it}.
\]
Here, $Y_{it}$ denotes annual wage and salary earnings; $C_i$ and $L_i$ are indicators
of college degree completion and lower family background. $\delta$ measures the
additional college premium for students from lower-background groups, $\lambda_t$
controls for year effects, and $X_{it}$ is a vector of the remaining control variables.

I estimate weighted pooled OLS with year effects, applying the same custom person
weight in each year, and cluster standard errors by person.

\section*{RESULTS}
% These are the previous report's exploratory results, retained as the revision
% baseline. Update the prose after changing the sample or model specification.
The weighted adjusted college earnings premium is \$25,225 for the lower-background
group and \$34,265 for the other group. The estimated difference is
$\hat\delta=-\$9,041$ ($p=0.140$; 95\% confidence interval: $-\$21,049$ to \$2,967).
The difference is not statistically significant and does not support my hypothesis
that lower-background students have a higher premium.\footnote{The cutoff was
selected after comparing several definitions and is therefore exploratory. The
60\% cutoff has the smallest $p$-value among the trials; nevertheless, the estimated
difference is statistically insignificant.}

% TODO: Choose the figure and regression-table presentation later.
\textit{[Insert your figure and its interpretation here.]}
\textit{[Insert your regression results table here.]}

Using the custom weights, 25.6\% of the lower-background group completed a bachelor's
degree or higher, compared with 58.6\% of the other group. This descriptive comparison
is consistent with my secondary hypothesis. These results describe weighted
associations within the complete-case sample. Standard errors do not account for
the full survey design. Missing data, topcoded earnings, and unobserved differences
limit the interpretation; the analysis does not establish a causal effect of college
completion on wealth inequality.

\section*{EXTENSIONS}
% TODO: Treat complementarity as an untested explanation, not an established
% finding; the interaction estimate is imprecise and the design is observational.
Based on my preliminary findings, the underlying mechanisms driving the observed
result remain uncertain. My revised mechanism hypothesis is that college is more
of a complement to family background than a substitute for it, which requires
further validation. This could explain why students from disadvantaged backgrounds
are less likely to finish their degrees and potentially have a relatively lower premium.

% TODO: Substantiate these channels. College selectivity is not defined for
% never-attenders and may be a pathway, so adding it as a control changes the question.
One additional theory for future research is that, even if students from disadvantaged
backgrounds indeed have a lower college premium, this might be because they tend to
attend less selective schools (affected by cost, nepotism, and unobserved academic
preparation). Hence, further research with additional control variables such as the
selectivity of the college attended can be helpful.

\section*{REFERENCES}
\begin{list}{}{\setlength{\leftmargin}{0.3in}
              \setlength{\itemindent}{-0.3in}
              \setlength{\itemsep}{0.6\baselineskip}
              \setlength{\parsep}{0pt}}
\item Becker, Gary S., and Nigel Tomes. 1986. ``Human Capital and the Rise and Fall of Families.'' \textit{Journal of Labor Economics} 4 (3, Part 2): S1--S39.
\item De Nardi, Mariacristina. 2015. ``Quantitative Models of Wealth Inequality: A Survey.'' NBER Working Paper 21106.
\item Hendricks, L., T. Koreshkova, and O. Leukhina. 2026. ``Expanding Access to Selective Colleges.'' Federal Reserve Bank of St. Louis Working Paper 2026-005.
\item Mincer, Jacob A. 1974. \textit{Schooling, Experience, and Earnings}. National Bureau of Economic Research.
\end{list}
\end{document}
